% 彗星（综述）
% license CCBYSA3
% type Wiki

本文根据 CC-BY-SA 协议转载翻译自维基百科\href{https://en.wikipedia.org/wiki/Comet}{相关文章}。

\begin{figure}[ht]
\centering
\includegraphics[width=8cm]{./figures/45ca94c34375617b.png}
\caption{} \label{fig_Comet_1}
\end{figure}
彗星是一类由冰组成的、体积较小的太阳系天体或星际天体。当其经过接近太阳的区域时，会因受热而开始释放气体，这一过程称为释气作用（outgassing）。这一过程会在彗核周围产生一个延展的、未受引力束缚的大气层，即彗发（coma），并有时形成由气体和尘埃组成的彗尾，这些物质自彗发中被吹离。上述现象源于太阳辐射及向外流动的太阳风等离子体对彗核施加的作用。彗核的大小从数百米到数十公里不等，由松散聚集的冰、尘埃以及小型岩石颗粒构成。彗发的尺度可达地球直径的 15 倍，而彗尾的长度甚至可能超过一个天文单位。如果距离足够近且亮度足够高，彗星无需借助望远镜即可从地球上观测到，并可能在天空中呈现出长达 30 度（相当于 60 个满月）的角距。自古以来，彗星已被许多文化和宗教所观察和记录。

彗星通常具有高离心率的椭圆轨道，其轨道周期范围广泛，从数年到可能长达数百万年。短周期彗星来源于柯伊伯带或其相关的散射盘，这些区域位于海王星轨道之外。长周期彗星则被认为起源于奥尔特云，一个由冰质天体组成的球形云层，其范围从柯伊伯带外侧延伸至朝向最近恒星方向一半的距离。长周期彗星通常由经过的恒星引力扰动或银河潮汐触发，开始向太阳方向运动。双曲线轨道彗星可能仅一次穿越内太阳系，然后被抛向星际空间。彗星的出现称为一次显现（apparition）。

多次接近太阳的灭绝彗星（extinct comets）会失去几乎全部挥发性冰与尘埃，并可能逐渐呈现出类似小型小行星的外观。一般认为，小行星与彗星的起源不同：小行星形成于木星轨道以内，而非太阳系外部。然而，主带彗星（main-belt comets）以及活跃的半人马族（centaurs）小天体的发现，使小行星与彗星之间的界限变得模糊。21 世纪初，一些具备长周期彗星轨道、但表现出内太阳系小行星特征的小天体被称为曼岛彗星（Manx comets）。它们仍被归类为彗星，例如 C/2014 S3（PANSTARRS）。在 2013 年至 2017 年间，共发现了 27 颗曼岛彗星。

截至 2021 年 11 月，共已知 4,584 颗彗星。然而，这仅占可能存在的彗星总量的一小部分，因为太阳系外部（奥尔特云）中类彗星天体的储库估计约为一万亿个。平均每年约有一颗彗星可被肉眼看见，尽管其中许多较为暗淡，外观并不特别醒目。特别明亮的彗星被称为“伟大彗星”（great comets）。人类已利用无人探测器访问过彗星，例如美国 NASA 的“深度撞击号”（Deep Impact），其通过撞击坦普尔 1 号彗星形成一个新陨坑以研究彗星内部；又如欧洲航天局（ESA）的“罗塞塔号”（Rosetta），其成为首个在彗星上着陆的机器人探测器。
\subsection{词源}
\begin{figure}[ht]
\centering
\includegraphics[width=8cm]{./figures/6a3173d4cf8911ea.png}
\caption{公元 729 年的《盎格鲁－撒克逊编年史》以及尊敬的比德（the Venerable Bede）的著作中均提到了彗星。} \label{fig_Comet_2}
\end{figure}
单词“comet（彗星）”源自古英语 cometa，其来源为拉丁语 comēta 或 comētēs，而这些拉丁词又是对希腊语 κομήτης 的罗马化形式，意为“留着长发的（事物）”。《牛津英语词典》指出，在希腊语中术语（ἀστὴρ） κομήτης 已经表示“长发之星，即彗星”。κομήτης 来源于动词 κομᾶν（koman），意为“留长发”，而该动词又派生自名词 κόμη（komē），意为“头发”，此词也被用来表示“彗尾”。

彗星的天文学符号（在 Unicode 中表示）为 U+2604 ☄ COMET，由一个小圆盘与三条类似头发的延伸构成。
\subsection{物理特征}
\begin{figure}[ht]
\centering
\includegraphics[width=8cm]{./figures/37fee7dd52548b31.png}
\caption{彗星的结构} \label{fig_Comet_3}
\end{figure}
\subsubsection{彗核}
\begin{figure}[ht]
\centering
\includegraphics[width=8cm]{./figures/43510b130873771c.png}
\caption{在航天器飞越期间拍摄的 103P/Hartley 彗核图像。该彗核长度约为 2 千米。} \label{fig_Comet_4}
\end{figure}
彗星的固体核心结构被称为彗核。彗核由岩石、尘埃、水冰以及冻结的二氧化碳、一氧化碳、甲烷和氨等物质的混合体构成。因此，根据弗雷德·惠普尔（Fred Whipple）的模型，它们常被形象地称为“肮脏的雪球”（dirty snowballs）。尘埃含量更高的彗星则被称为“冰质尘球”（icy dirtballs）。术语“冰质尘球”在 2005 年 7 月 NASA 的“深度撞击号”任务向坦普尔 1 号彗星发送“撞击器”探测器并观测到碰撞后开始使用。2014 年的研究表明，彗星类似于“炸冰淇淋”，即其表面由密集的晶体冰与有机化合物混合构成，而内部的冰则更冷且密度更低。

彗核表面通常干燥、多尘或岩质，表明冰层隐藏在厚度数米的表层壳体之下。彗核中含有多种有机化合物，其中可能包括甲醇、氰化氢、甲醛、乙醇、乙烷以及可能更复杂的分子，例如长链烃类与氨基酸。2009 年确认在 NASA “星尘号”（Stardust）任务回收的彗星尘埃中发现了氨基酸甘氨酸。2011 年 8 月，一份基于 NASA 对地球陨石研究的报告发表，提出 DNA 与 RNA 的组成单元（腺嘌呤、鸟嘌呤及相关有机分子）可能形成于小行星与彗星之中。

彗核外层表面具有极低反照率，使其成为太阳系中反射率最低的物体之一。贾托号（Giotto）探测器发现，哈雷彗星（1P/Halley）的彗核仅反射约四个百分点的入射光；而“深空一号”（Deep Space 1）发现波罗雷彗星（Comet Borrelly）的表面反射率甚至不足 3.0\%。作为比较，沥青的反射率为 7\%。彗核表面的暗色物质可能由复杂的有机化合物构成。太阳加热会驱散较轻、易挥发的化合物，留下更大的有机物，它们通常非常黑暗，如焦油或原油。彗核表面的低反射率使其能够吸收驱动释气过程的热量。

已观测到的彗核半径可达 30 千米，但要准确测定其大小并不容易。322P/SOHO 的彗核直径可能只有 100–200 米。尽管观测仪器的灵敏度不断提高，但仍缺乏对更小彗星的探测，这使一些研究者认为直径小于 100 米的彗星可能真实地数量稀少。已知彗星的平均密度估计约为 $0.6 g/cm^3$。由于质量较低，彗核不会在自引力作用下形成球形，因此呈现不规则形状。
\begin{figure}[ht]
\centering
\includegraphics[width=8cm]{./figures/7e9645911c29b572.png}
\caption{} \label{fig_Comet_5}
\end{figure}


